% Part: first-order-logic
% Chapter: model-theory
% Section: nonstandard-arithmetic

\documentclass[../../../include/open-logic-section]{subfiles}

\begin{document}

\olfileid{mod}{bas}{nsa}
\section{د حساب نامعياري مدلونه}

\begin{defn}
د حساب !!{language} دې $\Lang{L_N}$ وي، چې يو !!{constant} $\Obj{0}$، دوه‌ځايي !!{predicate} $<$، يوځايي !!{function} $\prime$، او دوه‌ځايي !!{function}s $+$ او $\times$ لري۔
\begin{enumerate}
\item د حساب \emph{معياري مدل} د $\Lang{L_N}$ هغه !!{structure} $\Struct{N}$ دے چې ساحه ئې $\Nat = \{ 0, 1, 2, \dots\}$ ده؛ $\Obj{0}$ د $0$ په توګه، $<$ پر~$\Nat$ د کوچنيوالي د اړيکې په توګه، او $\prime$، $+$ او $\times$ په ترتيب سره پر $\Nat$ د جانشين، جمع او ضرب په توګه تعبيروي۔
\item \emph{رښتينے حساب} تيوري $\Theory{N}$ ده۔
\end{enumerate}
\end{defn}

په $\Lang{L_N}$ کښې د کار پر مهال د $\Obj{0}^{\prime\dots\prime}$ په بڼه هر ترم، چې د جانشين تابع ئې پر $\Obj{0}$ باندې $n$ ځله لګېدلې وي، په~$\num{n}$ لنډوو۔

\begin{defn}
د $\Lang{L_N}$ يو !!{structure}~$\Struct{M}$ هله او يوازې هله \emph{معياري} دے چې $\Struct{N} \iso \Struct{M}$۔
\end{defn}

\begin{thm}\ollabel{thm:non-std}
د رښتيني حساب نامعياري !!{enumerable} مدلونه شته۔
\end{thm}

\begin{proof}
د نوي !!{constant}~$c$ په زياتولو $\Lang{L_N}$ غځوو او دا تيوري ګورو:
\[
\Theory{N} \cup \Setabs{\num{n} < c}{n \in \Nat}.
\]
تيوري په متناهي ډول د صدق وړ ده؛ نو د کمپکتوالي له مخې يو مدل $\Struct{M}$ لري، چې د ښکته لوېنهايم--سکولم قضيې له مخې ئې !!{enumerable} اخيستلے شو۔ که $\Domain{M}$ د $\Struct{M}$ ساحه وي، نو د $\Lang{L}$ د غيرمنطقي نښو تعبيرونه دې په $\Struct{M}$ کښې داسې وي: $\mathbf{z} = \Assign{\Obj{0}}{M} \in \Domain{M}$، ${\prec} =
\Assign{<}{M} \subseteq M^2$، $* = \Assign{\prime}{M} \colon
\Domain{M} \to \Domain{M}$، او $\oplus = \Assign{+}{M}, \otimes =
\Assign{\times}{M} \colon \Domain{M}^2 \to \Domain{M}$۔ د هر $x
\in \Domain{M}$ دپاره $x^*$ د $\Domain{M}$ هغه عنصر ښيي چې له $x$ څخه د~$*$ په لګولو ترلاسه کېږي۔

اوس، که $h$ د $\Struct{N}$ او د $\Struct{M}$ د اصلي حسابي راکمونې تر منځ ايزومورفيزم وے، نو داسې $n \in \Nat$ به وے چې $h(n) = \Assign{c}{M}$۔ نو په $\Struct{N}$ کښې داسې هره ارزښت ټاکنه $s$ واخلو چې $s(x) = n$۔ بيا $\Sat{N}{\eq[\num{n}][x]}[s]$؛ د \olref[iso]{thm:isom} د ثبوت له مخې $\Sat{M}{\eq[\num{n}][x]}[h \circ s]$ هم، نو $\Assign{c}{M} =
\mathbf{z}^{*\cdots *}$ (چې $*$ پکښې $n$ ځله تکرار شوے دے)۔ خو دا ناشوني دي، ځکه د فرض له مخې $\Sat{M}{\num{n} < c}$ او $\prec$ غيرانعکاسي ده۔ نو $\Struct{M}$ نامعياري دے۔
\end{proof}

\begin{prob}
پر سټ $X$ يوه اړيکه $R$ هله او يوازې هله \emph{ښه‌بنسټه} ده چې په~$R$ کښې نامتناهي ښکته تلونکې لړۍ نۀ وي؛ يعنې په $X$ کښې داسې $x_0$، $x_1$، $x_2$،~\dots نۀ وي چې $\dots x_2Rx_1Rx_0$۔ د دې فرض لاندې چې د زېرملو--فرېنکل سټ تيوري~$ZF$ سازګاره ده، وښيئ چې د $ZF$ نۀ ښه‌بنسټه مدلونه شته؛ يعنې داسې مدلونه $\mathfrak{M}$ چې $\dots x_2 \in x_1 \in x_0$۔
\end{prob}

په \olref{thm:non-std} کښې نامعياري مدل $\Struct{M}$ له معياري مدل سره ابتدايي هم‌ارز دے؛ ځکه ئې څو خاصيتونه راايستلے شو۔ د دې برخې پاتې متن همدې کار ته ځانګړے دے، چې د $\Struct{M}$ د خاصيتونو له لارې د $\Theory{N}$ د !!{enumerable} نامعياري مدلونو دقيقه ځانګړونه ترلاسه کړو۔

\begin{enumerate}
\item د $\Domain{M}$ هېڅ غړے له خپل ځان څخه د $\prec$ له مخې کوچنے نۀ دے: جمله $\lforall[x][\lnot x < x]$ په $\Struct{N}$ کښې رښتينې ده، نو په $\Struct{M}$ کښې هم۔
\item په ورته استدلال ترلاسه کوو چې $\prec$ د $\Domain{M}$ يو \emph{خطي ترتيب} دے؛ يعنې پر $\Domain{M}$ يوه کلي، غيرانعکاسي او متعدي اړيکه ده۔
\item عنصر $\mathbf{z}$ د $\Domain{M}$ د $\prec$ له مخې تر ټولو کوچنے عنصر دے۔
\item د $\Domain{M}$ هر غړے د $\prec$ له مخې له خپل $*$-جانشين څخه کوچنے دے، او $x^*$ د $\Domain{M}$ هغه تر ټولو کوچنے غړے دے چې د $\prec$ له مخې له $x$ څخه لوے دے۔
\item $\Struct{M}$ د $\prec$ يوه ابتدائي ټوټه لري چې له $\Nat$ سره ايزومورف ده: $\mathbf{z}, \mathbf{z}^*, \mathbf{z}^{**}, \dots$۔ دې ته د $\Domain{M}$ \emph{معياري برخه} وايو۔ د $\Domain{M}$ هر بل غړے \emph{نامعياري} دے۔ د $\Domain{M}$ نامعياري غړي خامخا شته؛ کنه په \olref{thm:non-std} کښې د ثبوت تابع $h$ به ايزومورفيزم وے۔ $n, m, \dots$ د هغو !!{variable}s په توګه کاروو چې د $\Struct{M}$ پر همدې معياري برخې ګرځي۔
\item هر نامعياري عنصر له هر معياري عنصر څخه لوے دے؛ ځکه د هر $n \in \Nat$ دپاره
  \[
  \Sat{N}{\lforall[z][(\lnot(\eq[z][\Obj{0}] \lor
  \dots \lor \eq[z][\num{n}]) \lif \num{n} < z)]},
  \]
  نو که $z \in \Domain{M}$ له ټولو معياري عناصرو بېل وي، بايد له ټولو څخه \emph{لوے} وي۔
\item له $\mathbf{z}$ پرته د $\Domain{M}$ هر غړے د $\Domain{M}$ د يوه يکتا عنصر $*$-جانشين دے؛ دغه مخکښېنی په $^*x$ ښيو۔ که $x = y^*$ وي، نو که د $x$ او $y$ يو معياري وي، دواړه معياري دي (او که يو نامعياري وي، دواړه نامعياري دي)۔
\item پر $\Domain{M}$ د هم‌ارزۍ اړيکه $\approx$ داسې تعريفوو: $x \approx y$ هله او يوازې هله چې د کوم \emph{معياري} $n$ دپاره $x \oplus n = y$ يا $y \oplus n =x$۔ په نورو ټکو، $x \approx y$ هله او يوازې هله چې د $x$ او $y$ تر منځ واټن متناهي وي۔ که $n$ او $m$ معياري وي، نو $n \approx m$۔ د $x$ \emph{بلاک} د هم‌ارزۍ ټولګے $[x] = \Setabs{y}{x
  \approx y}$ تعريفوو۔
\item فرض کړئ چې $x \prec y$ او $x \not\approx y$۔ څرنګه چې $\Sat{N}{\lforall[x][\lforall[y][(x < y \lif (x' < y \lor x' =
      y))]]}$، نو $x^* \prec y$ يا $x^* = y$۔ دويم حالت ناشونے دے، ځکه ترې $x \approx y$ راځي؛ نو $x^* \prec y$۔ \emph{OLCMP-108: سرچينه دلته پخوانے کمزورے شرط بيا تکراروي؛ د جانشين د مساوات له ردولو وروسته اړينه پايله د جانشين سخت کوچنيوالے دے۔ همدا پاتې اثباتي ګام څرګند شو، نۀ دا چې پخوانے کمزورے شرط په خپله ناسم وے۔} همداسې، که $x \prec y$ او $x \not\approx y$، نو $x \prec {^*y}$۔ ځکه که $x \prec y$ او $x \not\approx y$، نو هر $w \approx x$ د $\prec$ له مخې له هر $v \approx y$ څخه کوچنے دے۔ هر \emph{نامعياري} بلاک $[x]$ داسې دوه‌خوا نامتناهي لړۍ جوړوي:
  \[
  \cdots \prec  {^{**}x} \prec {^*}x \prec x \prec x^* \prec x^{**}
  \prec \cdots
  \]
  دې ته د $Z$ لړۍ وايو، ځکه د صحيح عددونو ترتيبي ډول لري۔ \emph{OLCMP-106: د سرچينې «هر بلاک» دلته نامعياري بلاک ته محدود شو۔ معياري بلاک تر ټولو کوچنے عنصر صفر لري او د طبيعي عددونو په شان يوې خوا ته نامتناهي دے؛ دوه‌خوا د صحيح عددونو لړۍ نۀ ده۔}
\item د $\prec$ ترتيب بلاکونو ته پورته کېدے شي: که $x \prec y$ وي او دواړه په بېلو بلاکونو کښې وي، نو د $x$ بلاک د $y$ له بلاک څخه کوچنے دے۔ \emph{OLCMP-107: سرچينه د بېلو بلاکونو شرط نۀ ليکي۔ بې له دې شرطه سخت ترتيب نۀ ترلاسه کېږي؛ صفر او يو بېل مرتب عناصر دي، خو په يوه معياري بلاک کښې دي۔} بلاک هله \emph{نامعياري} دے چې نامعياري عنصر ولري۔ معياري بلاک تر ټولو کوچنے بلاک دے۔
\item تر ټولو کوچنے نامعياري بلاک نشته: که $y$ نامعياري وي، نو داسې $x \prec y$ شته چې $x$ هم نامعياري وي او $x \not\approx y$۔ ثبوت: په معياري مدل $\Struct{N}$ کښې هر عدد پر دوو وېشل کېږي، او ښايي يو پاتې شي؛ يعنې
  $\Sat{N}{\lforall[y][\lexists[x][(y = x + x \lor y = x + x +
        \Obj{0}')]]}$۔ \emph{OLCMP-105: په سرچينه کښې د نيمې برخې شاهد دپاره هم عمومي کمي کوونکے ليکل شوے دے۔ دا ناسم دے؛ د صفر دپاره يو د نيمې برخې شاهد نۀ دے۔ دلته داخلي کمي کوونکے وجودي شو، چې د هر عدد دپاره د يوه مناسب شاهد شتون ووايي۔} د ابتدايي هم‌ارزۍ له مخې د هر $y \in \Domain{M}$ دپاره داسې $x \in \Domain{M}$ شته چې $x \oplus x = y$ يا $x \oplus x \oplus \mathbf{z}^*= y$۔ که $x$ معياري وے، $y$ به هم معياري وے؛ نو $x$ نامعياري دے۔ سربېره پر دې، $x$ او $y$ په بېلو بلاکونو کښې دي؛ يعنې $x \not\approx y$۔ د دې د ليدلو دپاره فرض کړئ چې دواړه په يوه بلاک کښې وي؛ يعنې د کوم معياري $n$ دپاره $x \oplus n = y$۔ که $y = x \oplus x$ وي، نو $x \oplus n = x
  \oplus x$، چې د جمع د حذف قانون له مخې ترې $x = n$ راځي (دغه قانون په $\Struct{N}$ کښې صدق کوي، نو په $\Struct{M}$ کښې هم)؛ بيا خو $x$ معياري شو۔ که $y = x \oplus x \oplus \mathbf{z}^*$ وي، استدلال همداسې دے۔
\item په ورته استدلال تر ټولو لوے بلاک هم نشته۔
\item د بلاکونو ترتيب ګڼ دے: که $[x]$ له $[y]$ څخه کوچنے وي (چې $x \not\approx y$)، نو د دواړو تر منځ يو له دواړو بېل بلاک $[z]$ شته۔ فرض کړئ چې $x \prec y$۔ لکه مخکې، $x \oplus y$ پر دوو وېشل کېږي (ښايي پاتې شونے ولري)؛ نو داسې $u \in \Domain{M}$ شته چې $x \oplus y = u \oplus u$ يا $x
  \oplus y = u \oplus u \oplus \mathbf{z}^*$۔ عنصر $u$ د $x$ او $y$ اوسط دے، نو د دواړو تر منځ دے۔ فرض کړئ چې $x \oplus y = u \oplus u$ (بل حالت ورته دے): که $u \approx x$ وي، نو د کوم معياري $n$ دپاره
  \[
  x \oplus y = x \oplus n \oplus x \oplus n,
  \]
  نو $y = x \oplus n \oplus n$، او $x \approx y$ به ولرو، چې له فرض سره خلاف دے۔ پايله دا ده چې $u \not\approx x$۔ په ورته استدلال $u \not\approx y$ هم ترلاسه کېږي۔
\end{enumerate}
نو نامعياري بلاکونه د ناطقو عددونو په شان مرتب دي: يو !!a{enumerable} خطي ترتيب جوړوي چې سرحدي ټکي نۀ لري، او ګڼوالے ئې هم ثابت شو۔ ترې راځي چې د رښتيني حساب د هر دوو !!{enumerable} نامعياري مدلونو $\Struct{M}_1$ او $\Struct{M_2}$ راکمونې هغې ژبې ته چې يوازې $<$ او $=$ لري، ايزومورف دي۔ په رښتيا، ايزومورفيزم $h$ داسې ټاکلے شو: د $\Struct{M_1}$ او $\Struct{M_2}$ معياري برخې له معياري مدل $\Struct{N}$ سره ايزومورف دي، نو له يو بل سره هم۔ د نامعياري برخې بلاکونه پخپله د ناطقو عددونو په شان مرتب دي، نو د \olref[dlo]{thm:cantorQ} له مخې ايزومورف دي؛ د بلاکونو ايزومورفيزم د بلاکونو \emph{دننه} ايزومورفيزم ته غځولے شو: په هرې جوړې کښې په خپل سر ټاکل شوي عناصر يو له بل سره برابر کړو، او بيا په $\Struct{M_1}$ کښې د $x$ د جانشين انځور په $\Struct{M_2}$ کښې د $x$ د انځور جانشين وټاکو۔ پام وکړئ چې ترې دا \emph{نۀ} راځي چې $\mathfrak{M}_1$ او $\mathfrak{M}_2$ د حساب په بشپړه ژبه کښې ايزومورف دي (ايزومورفيزم تل له يوې ټاکلې نښلنډې سره تړلے دے)، ځکه پر $\Domain{M_1}$ او $\Domain{M_2}$ د جمع او ضرب د تعريفولو نۀ ايزومورف لارې شته۔ (دا د واوت له مشهورې قضيې څخه هم راځي چې د يوې بشپړې تيورۍ د شمېرېدونکو مدلونو شمېر دوه نۀ شي کېدے۔)

\begin{prob}
وښيئ چې د حساب په نامعياري مدل کښې تر ټولو لوے بلاک نۀ شي کېدے۔
\end{prob}

\begin{prob}
د لومړۍ درجې !!{language} $\Lang{L}$ دې (له $\eq$ پرته) يوازې !!{predicate} $<$ ولري، او $\Struct{N} = (\Nat,
<)$ دې وي۔ د $\Struct{N}$ ټول متناهي يا هم‌متناهي فرعي سټونه د تعريف وړ دي۔ وښيئ چې د $\Struct{N}$ \emph{يوازې} همدا فرعي سټونه د تعريف وړ دي۔

(اشاره: لومړے، $\Obj{prc}(x,y)$ دې د $\Lang{L}$ هغه فارمول لنډ کړي چې وايي «$x$ د $y$ بې‌واسطې مخکښېنی دے»:
\[
x<y \land \lnot \lexists[z][(x<z \land z < y)].
\]
اوس د $\Struct{N}$ د هر تعريفېدونکي فرعي سټ سره په $\Lang{L}$ کښې يو فارمول $!A(x)$ تړلے دے۔ د هر داسې $!A$ دپاره دا جمله $!D$ وګورئ:
\[
\lexists[x][\lforall[y][\lforall[z][((x<y \land x<z \land \Obj{prc}(y,z)
  \land !A(y)) \lif !A(z))]]].
\]
وښيئ چې $\Sat{N}{!D}$ هله او يوازې هله چې د $\Struct{N}$ هغه فرعي سټ چې $!A$ ئې تعريفوي، متناهي يا هم‌متناهي وي۔

اوس $\Struct{M}$ دې له $\Struct{N}$ سره ابتدايي هم‌ارز نامعياري مدل وي۔ که $a \in \Domain{M}$ نامعياري وي، نو $b, c \in \Domain{M}$ دې له $a$ څخه لوے وي، او $b$ دې د~$c$ بې‌واسطې مخکښېنی وي۔ بيا د $\Domain{M}$ د ترتيبي جوړښت يو اتومورفيزم $h$ شته چې $h(b)=c$ (ولې؟)۔ ځکه که $b$ د $!A$ د صدق وړ کړي، نو $c$ هم همداسې کوي (ولې؟)۔ ترې راځي چې $!D$ په $\Struct{M}$ کښې رښتينې ده، نو په $\Struct{N}$ کښې هم۔ خو دا وايي چې د $\Struct{N}$ هغه فرعي سټ چې $!A$ ئې تعريفوي، متناهي يا هم‌متناهي دے۔
\end{prob}

\end{document}
